% 343_Zeta_Galois_FFGFT_De_ch.tex
% Johann Pascher, 2. September 2026
% Anlass: Fortsetzung von Dok. 342 — Konsequenzen der Galois-Faktorisierung
%         für die Zeta-Funktion des T⁴/Z₃-Torus

\providecommand{\BM}{\textbf{[B]}}
\providecommand{\KM}{\textbf{[K]}}
\providecommand{\SM}{\textbf{[S]}}

\chapter*{Die Zeta-Funktion des T$^4$/Z$_3$-Torus}
\addcontentsline{toc}{chapter}{Die Zeta-Funktion des T$^4$/Z$_3$-Torus}

\begin{abstract}
Die Galois-Faktorisierungsklassen aus Dok.~342 und die harmonische Reihe
auf dem T$^4$-Torus aus Dok.~159 führen gemeinsam zu einer präzisen Aussage
über die Zeta-Funktion der FFGFT-Raumzeit.

Vier Ergebnisse:
\textbf{(A)}~Die Z$_3$-Orbifold-Projektion von T$^4$/Z$_3$ transformiert
die volle Riemann-Zeta in eine Dirichlet-L-Funktion:
$\zeta_{T^4/\mathbb{Z}_3}(s) = L(s,\chi_0) = (1-3^{-s})\,\zeta(s)$ \BM.
\textbf{(B)}~Das Euler-Produkt dieser L-Funktion ist nach der
Eintrittstiefe $k = \mathrm{ord}_p(3)$ hierarchisch geordnet; die
physikalisch sichtbaren Primzahlen $\{2,3,5,7,11,13\}$ sind genau jene
mit $k \leq 6$ \BM.
\textbf{(C)}~Die Orbifold-Projektion erzeugt zusätzliche Nullstellen bei
$s = 2\pi i k/\ln 3$ ($k \in \mathbb{Z} \setminus \{0\}$) auf der
imaginären Achse $\mathrm{Re}(s) = 0$; diese sind die algebraische
Signatur der Z$_3$-Sektortrennung \BM.
\textbf{(D)}~Die vier Frobenius-Konjugationsklassen primitiver Elemente
in GF$(27)^*$ (Dok.~342, Satz~A) werden durch die FFGFT-Sektorpaarung
$k \leftrightarrow -k$ (Dok.~336; auf GF$(27)^*$: $x \mapsto x^{-1}$) zu
\emph{zwei} Klassen gepaart, $f_1 \leftrightarrow f_2$ und $f_3 \leftrightarrow f_4$;
die Negation $x \mapsto -x = x^{13}\,x$ koppelt jede primitive Klasse an genau
eine Ordnung-13-Klasse \BM. Der Galois-Charakter von
$\mathrm{Gal}(\mathrm{GF}(27)/\mathrm{GF}(3)) \cong \mathbb{Z}_3$ unterscheidet
die Klassen nicht; die Unterscheidung liegt in $(\mathbb{Z}/26)^*/\langle 3\rangle \cong \mathbb{Z}_4$.
Sektorpaarung und Majorana-Kriterium $k\mapsto k^{-1}$ (Dok.~340) zerlegen die
beiden Schichten deckungsgleich: $\{f_3,f_4\}$ Dirac (geladene Leptonen),
$\{f_1,f_2\}$ Majorana (Neutrinos; $f_1=O_1$ offen, Profil steriler Neutrinos) \BM/\SM.

Zwei weitere Ergebnisse:
\textbf{(E)}~Die harmonischen Primen $\{5,7,11,13\}$ sind durch die
Primfaktorzerlegungen von $3^k-1$ für $k=3,4,5,6$ eindeutig erzwungen;
$p=13$ ist die weltweit einzige Primzahl mit $\mathrm{ord}_p(3)=3$ \BM.
\textbf{(F)}~Die Euler-Faktor-Schwelle bei $p^*\approx9{,}76$ ($s=2$) markiert
das 7-limit exakt; ab $p_\mathrm{max}\approx30$ liegt die Galois-Produktdichte
auf Zufallsniveau; die Weyl-Obstruktion (Dok.~316) zeigt, dass die Grenze
nicht im Werkzeug liegt; Shors Schaltkreismodell ist digital und davon nicht betroffen, der physische
Apparat aber ein Resonanzsystem, das den Resonanzgrenzen unterliegt, solange das
Schwellentheorem bei Skalierung nicht demonstriert ist \SM; empirisch ist kein
Quantenvorteil nachgewiesen [X].

Prüfskripte: \texttt{pruef\_343\_zeta\_galois.py},
\texttt{pruef\_343b\_klassen\_symmetrie.py},
\texttt{pruef\_344\_primzahlmuster.py}.
\end{abstract}

% ============================================================
\section{Statuszeichen und Ausgangslage}
% ============================================================

Die folgenden Statusmarkierungen werden durchgängig verwendet:
\begin{center}
\begin{tabular}{@{}lp{10.5cm}@{}}
\textbf{[K]} & \emph{Konfirmiert} — vollständig aus $\xi$ oder einer korpusgesicherten
  Ableitung bewiesen; alle Assertions im Prüfskript bestanden. \\[4pt]
\textbf{[B]} & \emph{Belegt} — numerisch gesichert; analytischer Beweis vorhanden oder trivial. \\[4pt]
\textbf{[S]} & \emph{Spekulativ} — Hinweis oder Vermutung, noch nicht vollständig bewiesen. \\
\end{tabular}
\end{center}

Dok.~159 zeigt: Die spektrale Zeta-Funktion des Laplace-Operators auf einem
Torus $T^d$ ist die Epstein-Zeta-Funktion, deren eindimensionaler Spezialfall die
Riemann-Zeta $\zeta(s) = \sum_{n=1}^\infty n^{-s}$ ist.
Dok.~342 zeigt: Die Primzahlen erscheinen in der GF$(3^k)$-Hierarchie des
T$^4$/Z$_3$-Orbifolds nach der Eintrittstiefe $k = \mathrm{ord}_p(3)$ geordnet.

Die vorliegende Arbeit verbindet beide Ergebnisse: Was bedeutet die
GF$(3^k)$-Hierarchie für die Zeta-Funktion selbst?

% ============================================================
\section{Satz A: Orbifold-Projektion und Dirichlet-L-Funktion}
% ============================================================

\subsection{Zerlegung der Zeta nach Z$_3$-Klassen}

Auf T$^4$/Z$_3$ sind die erlaubten Moden $n \in \mathbb{Z}^4$ durch die
Z$_3$-Invarianz strukturiert. Im eindimensionalen Analogon zerfällt
$\zeta(s)$ in drei Teilsummen nach $n \bmod 3$ \BM:

\begin{equation}
  \zeta(s) = \underbrace{\sum_{\substack{n=1\\3\mid n}}^\infty \frac{1}{n^s}}_{= 3^{-s}\zeta(s)}
  + \underbrace{\sum_{\substack{n=1\\n\equiv 1(3)}}^\infty \frac{1}{n^s}}_{L_1(s)}
  + \underbrace{\sum_{\substack{n=1\\n\equiv 2(3)}}^\infty \frac{1}{n^s}}_{L_2(s)}
\end{equation}

Der Fixpunkt-Sektor ($n \equiv 0$) liefert genau $3^{-s}\zeta(s)$; die
Twist-Sektoren ($n \equiv \pm 1$) zusammen $(1 - 3^{-s})\zeta(s)$.

\subsection{Hauptcharakter $\chi_0$ mod~3}

Der Dirichlet-Charakter $\chi_0 \bmod 3$ mit $\chi_0(n) = 0$ falls $3 \mid n$,
sonst $\chi_0(n) = 1$ liefert \BM:
\begin{equation}
  L(s,\chi_0) = \sum_{\gcd(n,3)=1} \frac{1}{n^s}
  = (1 - 3^{-s})\,\zeta(s)
  = \prod_{p \neq 3} \frac{1}{1-p^{-s}}
\end{equation}

Dies ist die \textbf{spektrale Zeta-Funktion des T$^4$/Z$_3$-Torus}: Die
Charakteristik $p=3$ tritt nicht als Euler-Faktor auf — sie ist die
geometrische Grundlage des Orbifolds und steht außerhalb der Primzahl-Hierarchie.

\subsection{Nicht-trivialer Charakter $\chi_1$ mod~3}

Der Charakter $\chi_1 \bmod 3$ mit $\chi_1(1) = +1$, $\chi_1(2) = -1$,
$\chi_1(0) = 0$ misst die \emph{Asymmetrie} der Twist-Sektoren \BM:
\begin{equation}
  L(s,\chi_1) = \sum_{\substack{n=1\\n\equiv 1(3)}}^\infty \frac{1}{n^s}
              - \sum_{\substack{n=1\\n\equiv 2(3)}}^\infty \frac{1}{n^s}
\end{equation}
mit dem exakten Wert (Dirichlet 1837):
\begin{equation}
  L(1,\chi_1) = \frac{\pi}{3\sqrt{3}}
\end{equation}

Physikalische Bedeutung \SM: $L(s,\chi_1)$ misst die Differenz zwischen
Wicklungen $w = +1$ und $w = -1$ auf dem Torus, also die $\mathbb{Z}_3$-Sektor\-asymmetrie.
Numerisch gilt $L(2,\chi_0) \approx 1.4622$ und $L(2,\chi_1) \approx 0.7813$.
Ob $L(s,\chi_1)$ eine direkte physikalische Observable in FFGFT ist
(Teilchen-Antiteilchen-Asymmetrie im Twisted Sector?), bleibt offen \SM.

% ============================================================
\section{Satz B: Hierarchisches Euler-Produkt nach $k = \mathrm{ord}_p(3)$}
% ============================================================

Die Primzahlen in $L(s,\chi_0) = \prod_{p\neq 3}(1-p^{-s})^{-1}$
erscheinen nicht gleichwertig. Nach der Eintrittstiefe $k = \mathrm{ord}_p(3)$
aus Dok.~342, Satz~B, ist das Euler-Produkt hierarchisch geordnet \BM:

\begin{equation}
  L(s,\chi_0) = \prod_{k=1}^{\infty} \prod_{\substack{p\neq 3\\ \mathrm{ord}_p(3)=k}}
  \frac{1}{1 - p^{-s}}
\end{equation}

\begin{table}[H]
\centering
\caption{Hierarchie der Euler-Faktoren nach $k = \mathrm{ord}_p(3)$}
\begin{tabular}{rrlll}
\toprule
$k$ & neue Primzahl(en) & Euler-Faktor bei $s=2$ & Harmonik & FFGFT-Rolle \\
\midrule
1 & 2 & $4/3 = 1.333$ & Oktave & Z$_2$-Fixpunkt, massiver Sektor \\
3 & 13 & $170/169 \approx 1.006$ & Tredezime & GF$(27)^*$, Leptonen \\
4 & 5 & $25/24 \approx 1.042$ & Große Terz & Faktor in $(m_\mu/m_e)^2$ \\
5 & 11 & $122/121 \approx 1.008$ & Undezime & Neutrinos (Dok.~340) \\
6 & 7 & $49/48 \approx 1.021$ & Septime & Top-Quark (Dok.~289) \\
$\geq 7$ & $\geq 41, 1093, \ldots$ & $\approx 1 + p^{-2}$ & keine kl. Harmonik & physikalisch unterdrückt \\
\bottomrule
\end{tabular}
\end{table}

\subsection{Fraktale Unterdrückung höherer Euler-Faktoren}

Im fraktalen T$^4$-Torus (Dok.~159) werden die Beiträge hoher Moden durch den
Dämpfungsfaktor $\tau^{-(1+\xi)k}$ unterdrückt. Dies überträgt sich auf die
Euler-Faktoren: Das \emph{gewichtete} Euler-Produkt lautet
\begin{equation}
  \zeta_{T^4/\mathbb{Z}_3}^{\,\mathrm{frak}}(s)
  \;\cong\; \prod_{k=1}^{\infty}\;
  \prod_{\substack{p\neq 3\\ \mathrm{ord}_p(3)=k}}
  \left(\frac{1}{1-p^{-s}}\right)^{\!\tau^{-k(1+\xi)}}
\end{equation}
mit $\xi = 4/30000$ und $\tau < 1$ dem fraktalen Dämpfungsparameter.
Da $\tau^{-k(1+\xi)} \to 0$ für $k \to \infty$, wird jeder Euler-Faktor
mit wachsendem $k$ auf~1 gedrückt: Primzahlen mit $\mathrm{ord}_p(3) \geq 7$
leisten keinen physikalischen Beitrag \BM.

Das ist dieselbe Aussage wie Dok.~342, Satz~D (Galois-Raster $<$ fraktale
Korrektur ab $k=4$), jetzt in der Sprache der Zeta-Funktion formuliert.

% ============================================================
\section{Satz C: Neue Nullstellen durch die Orbifold-Projektion}
% ============================================================

\subsection{Triviale Nullstellen von $(1-3^{-s})$}

Der Faktor $(1-3^{-s})$ in $L(s,\chi_0) = (1-3^{-s})\zeta(s)$
hat Nullstellen bei \BM:
\begin{equation}
  3^{-s} = 1 \quad\Longleftrightarrow\quad s = \frac{2\pi i\, k}{\ln 3},
  \qquad k \in \mathbb{Z} \setminus \{0\}
\end{equation}

\begin{table}[H]
\centering
\caption{Nullstellen der Orbifold-Projektion auf der imaginären Achse}
\begin{tabular}{rcc}
\toprule
$k$ & $s = 2\pi i k/\ln 3$ & $\mathrm{Im}(s)$ \\
\midrule
$\pm 1$ & $\pm 5{,}719\,i$ & $\pm 2\pi/\ln 3$ \\
$\pm 2$ & $\pm 11{,}438\,i$ & \\
$\pm 3$ & $\pm 17{,}157\,i$ & \\
\bottomrule
\end{tabular}
\end{table}

Diese Nullstellen liegen auf $\mathrm{Re}(s) = 0$ (imaginäre Achse), \emph{nicht}
auf der kritischen Geraden $\mathrm{Re}(s) = 1/2$ der Riemann-Hypothese.
Sie sind dem Spektrum der Riemann-Zeta fremd — sie entstehen ausschließlich
durch die Z$_3$-Orbifold-Projektion.

\subsection{Physikalische Bedeutung der Nullstellen}

In der Torus-Sprache (Dok.~159) entsprechen Nullstellen der spektralen
Zeta-Funktion Moden, bei denen die Wellenfunktion eine Knotenstruktur
auf dem Torus hat. Die neuen Nullstellen bei $s = 2\pi i k/\ln 3$
entsprechen genau jenen Moden, die der Z$_3$-Fixpunktbedingung
$n \equiv 0 \pmod{3}$ unterliegen und im Twist-Sektor-Spektrum fehlen.

Die Nullstellen-Abstände $2\pi/\ln 3 \approx 5{,}72$ kodieren die
Periodizität der Z$_3$-Gitterstruktur. Ein Resonanztest der ersten
acht bekannten Riemann-Nullstellen $\gamma_n$ gegen diese Einheit
zeigt keine signifikante Übereinstimmung (Abweichungen $0{,}15$--$0{,}47$) \BM:
Die volle Riemann-Zeta kodiert die Verteilung \emph{aller} Primzahlen,
nicht nur der Z$_3$-kompatiblen.

\subsection{Vollständige Nullstellenstruktur der Orbifold-Zeta}

$L(s,\chi_0)$ hat drei Typen von Nullstellen:
\begin{itemize}
  \item \textbf{Triviale Nullstellen von $\zeta(s)$}: bei $s = -2, -4, -6, \ldots$
        (negative gerade Ganzzahlen)
  \item \textbf{Nicht-triviale Riemann-Nullstellen}: auf $\mathrm{Re}(s) = 1/2$
        (angenommen), unverändert übernommen
  \item \textbf{Orbifold-Nullstellen}: bei $s = 2\pi i k/\ln 3$, $k \neq 0$
        — neu, durch Z$_3$-Projektion erzeugt \BM
\end{itemize}

% ============================================================
\section{Satz D: Die vier Klassen unter der Sektorpaarung}
% ============================================================

\subsection{Der Galois-Charakter unterscheidet die Klassen nicht}

Die Galois-Gruppe $\mathrm{Gal}(\mathrm{GF}(27)/\mathrm{GF}(3)) = \langle F \rangle \cong \mathbb{Z}_3$
mit $F\colon x \mapsto x^3$ hat drei Charaktere $\chi_j(F) = \omega^j$, $j = 0,1,2$.
Alle vier primitiven Klassen $f_1, \ldots, f_4$ aus Dok.~342 sind \emph{reguläre}
$F$-Orbits der Länge~3: $F$ wirkt auf jedem als Dreierzyklus, und der zugehörige
Charakter ist für alle vier derselbe. Der Galois-Charakter trennt die Klassen
also nicht \BM.

Die Unterscheidung liegt im diskreten Logarithmus: Mit einem Erzeuger $x$ von
GF$(27)^*$ ist jede Klasse ein Orbit $\{x^k, x^{3k}, x^{9k}\}$ mit $\gcd(k,26)=1$,
also eine Nebenklasse von $\langle 3\rangle = \{1,3,9\}$ in
$(\mathbb{Z}/26)^* \cong \mathbb{Z}_{12}$. Es gibt $12/3 = 4$ Nebenklassen \BM:

\begin{table}[H]
\centering
\caption{Die vier primitiven Klassen als Nebenklassen von $\langle 3\rangle$ in $(\mathbb{Z}/26)^*$
  (Basis $x^3 + 2x + 1 = 0$)}
\begin{tabular}{llll}
\toprule
Klasse & Minimalpolynom & Exponenten $k$ & Rolle \\
\midrule
$f_1$ & $x^3 + 2x + 1$ & $\{1, 3, 9\}$ & Majorana-Schicht (Satz~D$'''$) \\
$f_2$ & $x^3 + 2x^2 + 1$ & $\{17, 23, 25\}$ & $= f_1^{-1}$, Majorana-Schicht \\
$f_3$ & $x^3 + 2x^2 + x + 1$ & $\{5, 15, 19\}$ & Dirac-Schicht \\
$f_4$ & $x^3 + x^2 + 2x + 1$ & $\{7, 11, 21\}$ & $= f_3^{-1}$, Dirac-Schicht \\
\bottomrule
\end{tabular}
\end{table}

Die Faktorgruppe $(\mathbb{Z}/26)^*/\langle 3\rangle \cong \mathbb{Z}_4$ ist zyklisch;
die Inversion $k \mapsto -k$ ist ihr Element der Ordnung~2.

\subsection{Sektorpaarung reduziert vier auf zwei}

Dok.~336 identifiziert die FFGFT-Sektorpaarung $k \leftrightarrow -k$ mit dem
Frobenius-Automorphismus; auf der Ebene der Elemente von GF$(27)^*$ ist sie die
Inversion $x \mapsto x^{-1} = x^{25}$. Sie permutiert die Klassen als \BM
\begin{equation}
  f_1 \longleftrightarrow f_2, \qquad f_3 \longleftrightarrow f_4 .
\end{equation}
Ist die Sektorpaarung eine Symmetrie der Massenschicht --- im Korpus ist sie das ---,
so gibt es nicht vier, sondern \textbf{zwei} physikalisch unterscheidbare
GF$(27)^*$-Schichten: $\{f_1, f_2\}$ und $\{f_3, f_4\}$ (Zuordnung in Satz~D$'''$).

\subsection{Negation koppelt Ordnung 26 an Ordnung 13}

Wegen $-1 = x^{13}$ in GF$(27)$ bildet $x \mapsto -x$ jedes Element der Ordnung~26
auf ein Element der Ordnung~13 ab und umgekehrt. Die vier Ordnung-13-Klassen
$g_1, \ldots, g_4$ aus Dok.~342 sind daher keine eigenständige Struktur, sondern
die Vorzeichenpartner der primitiven Klassen \BM:
\begin{equation}
  f_i \leftrightarrow g_i \qquad (i = 1, \ldots, 4).
\end{equation}
Der in Dok.~342 als offen geführte Kandidat einer ,,Sub-Leptonen-Struktur``
aus den Ordnung-13-Polynomen entfällt damit.

\subsection{Offene Frage}

Die Reduktion $4 \to 2$ verschärft die Frage aus Dok.~342: Nicht ,,tragen
$f_2, f_3, f_4$ weitere Schichten``, sondern ,,es gibt genau zwei
GF$(27)^*$-Schichten $\{f_1, f_2\}$ und $\{f_3, f_4\}$ --- welche Teilchen
tragen welche?{}`` Satz~D$'''$ beantwortet das aus der Algebra.

\subsection{Satz D$'''$: Quantenzahlen der zweiten Schicht}

Welche Quantenzahlen hätten die beiden Schichten? Die Algebra beantwortet
das ohne physikalische Annahme (\texttt{pruef\_343f\_zweite\_schicht.py}). Jede
primitive Klasse $f_i$ liegt in genau einem $\mathbb{Z}_{13}$-Orbit des
Frobenius (Dok.~339); mit den Exponenten aus Tabelle~3: $f_1=\{1,3,9\} \mapsto O_1$,
$f_2=\{17,23,25\} \mapsto O_3=\{4,10,12\}$, $f_3=\{5,15,19\} \mapsto O_2=\{2,5,6\}$,
$f_4=\{7,11,21\} \mapsto O_4=\{7,8,11\}$. Dok.~340 nimmt die
Selbstinversheit eines Orbits unter $k \mapsto k^{-1} \pmod{13}$ als
Majorana-Kriterium. Beide Involutionen zusammen ergeben \BM:

\begin{table}[H]
\centering
\caption{Quantenzahlen der vier primitiven Klassen aus der Algebra}
\begin{tabular}{llllll}
\toprule
Klasse & $\mathbb{Z}_{13}$-Orbit & $k\mapsto -k$ & $k\mapsto k^{-1}$ & Profil & Schicht \\
\midrule
$f_1$ & $O_1$ & $\to f_2$ & selbstinvers & Majorana & $\{f_1,f_2\}$ \\
$f_2$ & $O_3$ & $\to f_1$ & selbstinvers & Majorana & $\{f_1,f_2\}$ \\
$f_3$ & $O_2$ & $\to f_4$ & $\to O_4$ & Dirac & $\{f_3,f_4\}$ \\
$f_4$ & $O_4$ & $\to f_3$ & $\to O_2$ & Dirac & $\{f_3,f_4\}$ \\
\bottomrule
\end{tabular}
\end{table}

Die Sektorpaarung $k \mapsto -k$ und das Majorana-Kriterium $k \mapsto k^{-1}$
zerlegen die vier Klassen deckungsgleich: \textbf{$\{f_1,f_2\}$ ist die
Majorana-Schicht, $\{f_3,f_4\}$ die Dirac-Schicht} \BM. Mit der Negation
$f_i \leftrightarrow g_i$ ist die Nummerierung aus Dok.~342 damit die natürliche:
Sektorpaarung paart benachbarte Indizes, Negation gleiche Indizes. Beide Schichten tragen
$\mathbb{Z}_2$-Parität~1 (massiver Sektor, Dok.~339), keine Farbe und drei
Zustände je Klasse (Dreier-Orbit).

\textbf{Korrektur.} Dok.~341 legt für die Leptonen nur die Ordnung~26 in
G$(6)$ fest, keine Klasse; frühere Zuschreibungen einer Leptonenklasse zu
Dok.~341 waren nicht gedeckt. Die Zuordnung folgt erst hier: Geladene Leptonen
sind Dirac-Teilchen und gehören in $\{f_3,f_4\}$.

\textbf{Deutung \SM.} Die zweite Schicht ist damit nicht neu, sondern der
Neutrinosektor: Dok.~340 setzt die aktiven Neutrinos in $O_3 = f_2$. Offen
bleibt $f_1 = O_1 = \{1,3,9\}$ --- der Orbit des Erzeugers selbst: Majorana,
drei Zustände, $\mathbb{Z}_2$-Parität~1, Sektorpartner der aktiven Neutrinos
unter $k \mapsto -k$. Das ist das Profil steriler (rechtshändiger) Neutrinos,
die Dok.~340 bereits als Majorana im Orbit-3-Sektor mit $m_\mathrm{min} =
4\,m_\nu = 18{,}2$~meV vorhersagt; die Paarung $O_1 \leftrightarrow O_3$ unter
der Sektorpaarung ist algebraisch die Seesaw-Paarung aktiv$\leftrightarrow$steril.
Eine Kuriosität ohne Status: $(\max O_1)^2 - 1 = 80 = |\mathrm{GF}(81)^*|$,
analog zu $(\max O_4)^2 - 1 = 120 = \Delta m^2_\mathrm{atm}/m_\nu^2$ aus
Dok.~340. Ob $O_1$ die Massenskala der sterilen Neutrinos trägt, ist offen.


% ============================================================
\section{Satz E: Primzahlmuster in der GF$(3^k)$-Hierarchie}
% ============================================================

\subsection{Die harmonischen Primen sind eindeutig erzwungen}

Eine Primzahl $p$ erfüllt $\mathrm{ord}_p(3) = k$ genau dann, wenn $p$
ein Primteiler von $3^k - 1$ ist und $p \nmid 3^j - 1$ für alle $j < k$.
Die Primfaktorzerlegungen der ersten relevanten Werte ergeben \BM:

\begin{table}[H]
\centering
\caption{Primfaktoren von $3^k - 1$ und die erzwungenen physikalischen Primzahlen}
\begin{tabular}{rrlll}
\toprule
$k$ & $3^k - 1$ & Primfaktoren & phys. Prime & Anmerkung \\
\midrule
1 & 2 & $\{2\}$ & 2 & Oktave \\
3 & 26 & $\{2, 13\}$ & \textbf{13} & einzige Primzahl mit $\mathrm{ord}_p(3)=3$ weltweit \\
4 & 80 & $\{2, 5\}$ & \textbf{5} & einziger Nicht-2-Primteiler \\
5 & 242 & $\{2, 11\}$ & \textbf{11} & einziger Nicht-2-Primteiler \\
6 & 728 & $\{2, 7, 13\}$ & \textbf{7} & (13 hat $\mathrm{ord}_{13}(3)=3$, nicht 6) \\
\bottomrule
\end{tabular}
\end{table}

Das stärkste Ergebnis betrifft $p = 13$ \BM: Da $3^3 - 1 = 26 = 2 \cdot 13$
und $\mathrm{ord}_2(3) = 1$, ist $p = 13$ die \emph{einzige} Primzahl
weltweit mit $\mathrm{ord}_p(3) = 3$. Es gibt keine weitere, und es kann keine
geben. Die physikalischen Primen $\{5, 7, 11, 13\}$ sind damit nicht durch
harmonische Auswahl ausgezeichnet, sondern durch die Primfaktorstruktur von
$3^k - 1$ für $k = 3, 4, 5, 6$ vollständig und eindeutig erzwungen.

\subsection{Notwendige Kongruenzbedingung}

Aus $\mathrm{ord}_p(3) = k$ mit $3 \mid k$ folgt $3 \mid (p-1)$,
also $p \equiv 1 \pmod{3}$ \BM. Beweis: $p \mid 3^k - 1$ und $3 \mid k$
implizieren $p \mid (3^{k/3})^3 - 1$; da $p \nmid 3^{k/3} - 1$ (nach
Definition der Ordnung), folgt aus dem Faktorisierungssatz für
$x^3 - 1 = (x-1)(x^2+x+1)$, dass $p \mid x^2 + x + 1$ für $x = 3^{k/3}$,
was $3 \mid p - 1$ erzwingt. Insbesondere sind $p = 13$ (bei $k=3$) und
$p = 7$ (bei $k=6$) beide $\equiv 1 \pmod{3}$ — genau die Z$_3$-kompatiblen
Primzahlen des Korpus (Dok.~342, Satz~B).

% ============================================================
\section{Satz F: Grenzen der Unterscheidbarkeit und des Rechenaufwands}
% ============================================================

\subsection{Exakte Schwelle: Euler-Faktor gegen fraktale Korrektur}

Der Euler-Faktor $1/(1-p^{-s})$ überschreitet die fraktale Korrektur
$1{,}05\,\%$ (Dok.~338, bare $\to$ gemessen) genau für $p < p^*$
mit $p^* = (0{,}0105)^{-1/s}$ \BM:

\begin{table}[H]
\centering
\caption{Exakte Schwelle $p^*$ nach Wert von $s$}
\begin{tabular}{rrlll}
\toprule
$s$ & $p^*$ & Primzahlen $\leq p^*$ & harmonisches Limit & Euler-Faktor $\geq 1{,}05\,\%$ \\
\midrule
1 & 95{,}2 & $\{2,3,5,\ldots,89\}$ & --- & alle kleinen Primen \\
2 & 9{,}76 & $\{2, 3, 5, 7\}$ & 7-limit & genau Oktave, Quinte, Terz, Septime \\
3 & 4{,}57 & $\{2, 3\}$ & 3-limit & nur Oktave und Quinte \\
\bottomrule
\end{tabular}
\end{table}

Für $s = 2$ (Referenzwert der spektralen Zeta-Funktion) beträgt die Schwelle
$p^* \approx 9{,}76$: Die Primzahlen $\{2, 5, 7\}$ haben einen Euler-Faktor,
der die fraktale Korrektur übersteigt; $p = 11$ und alle größeren
unterschreiten sie. Das ist exakt das harmonische 7-limit.
$p = 11$ erscheint im Korpus (Neutrinos, Dok.~340) als Strukturformel
$\Delta m^2 = (11^2-1)\,m_\nu^2$, nicht als Galois-Produkt-Treffer.

\subsection{Galois-Produkte gegen Zufallsprodukte: Trefferrate}

Bildet man alle Produkte aus bis zu drei Galois-Bausteinen
$\{3^k, 3^k-1, \text{Primteiler}\}$ für $p \leq p_\mathrm{max}$ mit Exponenten
$\{-1, 1, 2\}$, so wächst die Produktmenge schnell und die Trefferrate
für beliebige Ziele erreicht das Zufallsniveau \BM:

\begin{table}[H]
\centering
\caption{Trefferrate physikalischer Verhältnisse vs.\ zufälliger Zahlen bei $\pm 0{,}5\,\%$}
\begin{tabular}{rrrrr}
\toprule
$p_\mathrm{max}$ & Werte & $43200$ & $m_\mu/m_e$ & Zufalls-Ø \\
\midrule
13  & 1334  & 2 & 0 & 0{,}92 \\
30  & 5802  & 3 & 2 & 3{,}60 \\
50  & 8946  & 3 & 2 & 5{,}33 \\
100 & 13398 & 4 & 3 & 7{,}66 \\
300 & 24992 & 8 & 12 & 14{,}15 \\
\bottomrule
\end{tabular}
\end{table}

Ab $p_\mathrm{max} \approx 30$ liegt die Trefferrate auf Zufallsniveau
(±0,5\,\%); bei der strengen Schwelle ±0,05\,\% bereits ab $p_\mathrm{max} \approx 13$.
Die Identität $43200 = |\mathrm{GF}(9)^*|^2 \cdot 5^2 \cdot |\mathrm{GF}(27)|$
(Dok.~338) ergibt sich ausschließlich aus Bausteinen mit $p \leq 5$ und ist
\emph{geometrisch} aus $\xi$ hergeleitet, nicht durch Produktsuche gefunden.

\subsection{Strukturelle Rechengrenze: Weyl-Obstruktion und Quantencomputer}

Die Unterscheidbarkeitsgrenze bei $p_\mathrm{max} \approx 13$ ist eine
\emph{physikalische} Grenze (fraktale Korrektur). Unabhängig davon besteht
im Korpus eine schärfere \emph{mathematische} Grenze, die jede Technologie
betrifft \BM:

\begin{keyresult}[Weyl-Obstruktion, Dok.~316]
Es gibt keine kompakte Riemannsche Mannigfaltigkeit — gleich welcher
Dimension, Krümmung, Topologie oder Orbifold-Struktur —, deren
Laplace-Spektrum die Riemann-Nullstellen sind.
Grund: $T\ln T$ ist keine Potenz $T^{d/2}$.
\end{keyresult}

Aus Dok.~147 folgt, dass ein $n$-Qubit-Register ein kompakter Zustandsraum
mit Eigenwertdichte $N(\lambda) \sim \lambda^{d/2}$ ist; arithmetische
Zähldichten wachsen dagegen wie $T \ln T$. Ein endlicher Zustandsraum kann
ein solches Spektrum nicht vollständig tragen, \emph{unabhängig von der
Registergröße} \BM.

Dok.~075 benennt die Konsequenz für Shors Algorithmus:

\begin{center}
\begin{tabular}{@{}lll@{}}
\toprule
Verfahren & Grenze & Typ \\
\midrule
Probeteilung & $O(\sqrt{n})$ & Rechenaufwand \\
Pollard $\rho$ & $O(n^{1/4})$ & Rechenaufwand \\
Shors QFT & Ressourcengrenze & Weyl-Obstruktion \\
\bottomrule
\end{tabular}
\end{center}

Die zentrale Aussage aus Dok.~316 gilt technologieunabhängig:
\begin{quote}
\emph{Die Grenze liegt nicht im Werkzeug, sondern im Gegenstand.}
Transformation (Fourier, QFT, Optik) projiziert und erzeugt keine
Information. Relaxation findet, was im Potential steht. Beide setzen
voraus, dass die gesuchte Struktur bereits vorhanden ist.
\end{quote}

Damit stehen physikalische Grenze ($p^* \approx 9{,}76$, Satz~F.1) und
mathematische Grenze (Weyl-Obstruktion, Satz~F.3) im selben Bild:
Ab einem bestimmten Punkt ist die Primzahlstruktur im Rauschen verborgen,
und kein analoges Werkzeug holt sie heraus; Shors Algorithmus ist als Modell digital,
als Apparat ein Resonator --- und damit bis zum Beweis des Gegenteils ebenfalls gebunden (Satz~G$'$).

% ============================================================
\section{Im Klartext: Was die Grenze bedeutet}
% ============================================================

\textbf{Frage.} Es lassen sich also beliebig dichte Harmonische berechnen ---
aber ab einer bestimmten Tiefe nicht mehr unterscheiden?

\textbf{Antwort.} Ja, genau das.

\emph{Berechnen kann man beliebig weit.} GF$(3^k)$ für $k = 7, 8, 12, 100$ ---
die Algebra hört nicht auf. Jedes $k$ liefert neue Primzahlen, neue
Polynomklassen, neue Akkorde. Bei $k = 12$ sind alle vier harmonischen Primen
$5, 7, 11, 13$ gleichzeitig in einem Körper; bei $k = 7$ tritt die
Wieferich-Primzahl $1093$ ein. Diese Strukturen sind exakt berechenbar und
mathematisch alle vorhanden.

\emph{Aber ab einer bestimmten Tiefe sind sie physikalisch bedeutungslos},
und zwar aus zwei Gründen, die zusammenfallen.

Erstens wird das Raster zu fein. Bei $k \leq 3$ liegen die erreichbaren
Zahlenwerte etwa $3\,\%$ auseinander --- trifft $43200$ einen Massenwert auf
$1\,\%$, ist das ein Ergebnis. Bei $k \geq 5$ liegen sie unter $0{,}5\,\%$
auseinander. Dann trifft \emph{jede} Zahl irgendetwas, und ein Treffer sagt
nichts mehr. Der Kontrolltest in Satz~F zeigt es: zweihundert Zufallszahlen
treffen genauso oft wie die echten Massenverhältnisse.

Zweitens --- und das ist der physikalische Grund --- ist die fraktale Korrektur
von etwa $1\,\%$ (Dok.~338, bare $\to$ gemessen) selbst schon eine Aussage
über höhere Wicklungsmoden. Sobald der Galois-Raster feiner als diese
Korrektur ist, unterscheidet die Natur selbst nicht mehr zwischen benachbarten
Galois-Werten. Die Struktur ist nicht weg; sie ist im Rauschen der höheren
Moden aufgegangen.

Die Analogie zur Harmonik ist exakt: Die Obertöne $2{:}1$ und $3{:}2$ hört
jeder. $17{:}16$ ist ein Halbton, $33{:}32$ ein Viertelton. Ab etwa dem
250.\ Oberton liegt der Abstand unter 5~Cent --- mathematisch exakt vorhanden,
akustisch nicht mehr unterscheidbar. Die Reihe wird zum Kontinuum.

\emph{Die Konsequenz für FFGFT.} Die exakte Galois-Schicht reicht bis GF$(27)$,
also $k = 3$: die Leptonenschicht, $43200$, die vier Klassen, $p = 13$ als
einzige Primzahl mit $\mathrm{ord}_p(3) = 3$. Darüber ($k = 4$, GF$(81)$,
$p = 5$) ist bereits Übergang. Ab $k = 5$ ist alles, was man durch Suche findet,
nicht mehr von Zufall unterscheidbar --- und darf keinen Status erhalten, es sei
denn, es kommt direkt aus $\xi$.

Das ist kein Mangel der Theorie, sondern ihre Konsistenz: Dok.~159 sagt
Modenkonvergenz, Dok.~316 sagt Weyl-Obstruktion, Dok.~342 und dieses Dokument
sagen dasselbe algebraisch. Es gibt eine Grenze, sie liegt bei $k \approx 3$--$4$,
und sie kommt aus der Sache selbst, nicht aus unserem Rechenvermögen.

% ============================================================
\section{Im Klartext II: Synthese, Analyse, Faktorisierung --- und die Kopplung}
% ============================================================

\textbf{Frage.} Was bedeutet die Grenze für Synthese und Analyse --- und für
die Faktorisierung? Und wie wirkt sich die Kopplung aus: Dok.~328 kennt drei
Kopplungsregime; ist bei reinen Zahlen dasselbe zu sehen wie bei physikalischen
Vorgängen?

\subsection{Synthese ist unbegrenzt, Analyse ist begrenzt}

\emph{Synthese} --- Obertöne addieren, Primzahlen multiplizieren,
GF$(3^k)$-Akkorde aus Partitionen bilden --- geht beliebig weit und beliebig
dicht. Nichts hindert daran, den 250.\ Oberton mit dem 251.\ zu überlagern oder
zwei 300-stellige Primzahlen zu multiplizieren. Die Synthese kennt die Grenze
aus Satz~F nicht.

\emph{Analyse} --- welche Obertöne stecken im Klang, welche Primzahlen in der
Zahl, welche Galois-Klasse trägt die Masse --- stößt genau an diese Grenze.
Das ist keine Frage des Rechenaufwands, sondern der Auflösung: Um Oberton $n$
von $n{+}1$ zu trennen, braucht man ein Zeitfenster von etwa $n$ Perioden; um
die Periode $r$ einer Zahl zu finden, eine Frequenzauflösung
$\Delta f \approx f_0/n^2$ (Dok.~075, Präzisionsbarriere). Ab einer Dichte
liegt der Abstand unter dem Rauschen, und dann ist die Zerlegung nicht mehr
eindeutig --- nicht weil sie nicht existiert, sondern weil viele Zerlegungen
gleich gut passen. Das ist der Zufallsniveau-Befund aus Satz~F in anderer
Sprache.

\subsection{Faktorisierung ist Analyse}

Eine Zahl $N = p \cdot q$ zu faktorisieren heißt, die zwei Grundtöne aus einer
Schwebung zu lesen. Bei kleinen $p, q$ --- dem 7-limit-Bereich, $k \leq 3$ ---
geht das; das sind genau die Primzahlen, die im FFGFT-Spektrum überhaupt
sichtbar sind. RSA-Primzahlen mit 300 Stellen liegen bei
$k = \mathrm{ord}_p(3)$ von der Größenordnung $p$ selbst, also in einer Tiefe,
die den Rasterabstand aus Satz~F um Hunderte Größenordnungen unterschreitet.
Dok.~316 formuliert es: Fourier, QFT und Optik projizieren, sie erzeugen keine
Information; die Periode muss in den präparierten Zuständen schon stecken, und
diese Präparation (modulare Exponentiation, $95\,\%$ der Gatterkosten,
Dok.~147) ist genau das Problem, das man lösen wollte.

Daraus folgt zweierlei. Erstens: FFGFT bricht RSA nicht --- und das ist kein
Mangel, sondern ein konsistentes Ergebnis. Die Asymmetrie
Synthese-leicht/Analyse-schwer, von der RSA lebt, ist dieselbe Asymmetrie, die
die Galois-Schicht bei $k \approx 3$--$4$ enden lässt. Zweitens, für den
Harmonikabau: Ein Instrument kann Obertöne beliebig dicht \emph{erzeugen}; Ohr
und Analyse \emph{hören} nur bis etwa zum 7-limit trennscharf. Was darüber
liegt, wirkt als Klangfarbe, nicht als Ton --- wie die höheren Wicklungsmoden
in FFGFT als $1\,\%$-Korrektur wirken, nicht als eigene Massenschicht.

\textbf{Vorbehalt \SM.} Die Übertragung von Satz~F auf $\mathbb{Z}/N\mathbb{Z}$
ist eine Analogie. Dok.~075 hält fest, dass $\xi$ in $\mathbb{Z}/N\mathbb{Z}$
keine Funktion hat und die Periode rein arithmetisch durch
$\mathrm{kgV}(p-1, q-1)$ bestimmt ist. Die Analogie trägt; die fraktale
Korrektur als \emph{Mechanismus} der Faktorisierungsgrenze nicht.

\subsection{Kopplung: dieselbe Dreiteilung, nicht dasselbe Gesetz}

Dok.~328 klassifiziert Kopplungen relativ zu den Verlustraten:
unterkritisch ($\kappa < 1$, ein Maximum), kritisch ($\kappa = 1$, maximal
flach), überkritisch ($\kappa > 1$, Aufspaltung). Einrasten auf ein
Frequenzverhältnis $p/q$ geschieht in Arnold-Zungen der Breite
$\Delta\Omega(p/q) \propto K^q$.

Bei reinen Zahlen übernimmt die Toleranz $\varepsilon$ die Rolle des
Kopplungsgrades: Eine reelle Zahl ,,rastet`` auf $p/q$ ein, wenn
$|x - p/q| < \varepsilon$. Die Brüche mit Nenner $\leq Q$ überdecken das
Intervall mit Maß $\approx 2\varepsilon \cdot 3Q^2/\pi^2$; Überdeckung~1
liefert die kritische Nennerhöhe \BM
\begin{equation}
  Q_c = \frac{\pi}{\sqrt{6\,\varepsilon}} .
\end{equation}
Für $\varepsilon = 1{,}05\,\%$ (fraktale Korrektur) ist $Q_c \approx 12{,}5$ ---
dieselbe Größenordnung wie $p_\mathrm{max} \approx 13$ in Satz~F. Die
Galois-Version bestätigt es numerisch (\texttt{pruef\_343c}): Bei
$\varepsilon = 1{,}05\,\%$ ist der Galois-Raster mit $p_\mathrm{max} = 7$
unterkritisch (Überdeckung $0{,}52$), mit $p_\mathrm{max} = 13$ kritisch
($0{,}92$), ab $p_\mathrm{max} \approx 30$ überkritisch ($1{,}00$).
\textbf{Die Grenze $k \approx 3$--$4$ aus Satz~F ist der kritische Punkt
$\kappa = 1$ im Sinne von Dok.~328} \BM.

Die Dreiteilung ist also in Arithmetik und Physik dieselbe. \emph{Das Gesetz
dahinter ist es nicht} \BM: Arnold-Zungen fallen exponentiell in $q$
($\Delta\Omega \propto K^q$), Farey-Abstände nur polynomial ($1/q^2$). Bei
$K = 0{,}3$ und $q = 5$ ist die Zunge bereits um den Faktor $10^{-5}$ schmaler
als das arithmetische Fangintervall; bei $q = 13$ um $10^{-15}$. Physikalisches
Einrasten ist exponentiell selektiv, arithmetische Unterscheidbarkeit
polynomial. Die Physik ,,hört`` weniger Harmonische, als die Arithmetik
zulässt --- die Kante ist dieselbe, aber physikalisch schärfer.

\subsection{FFGFT ist dynamisch tief unterkritisch}

Mit $K_\mathrm{eff} = 2\pi\xi \approx 8{,}4 \times 10^{-4}$ (Dok.~328) ist die
breiteste nichttriviale Arnold-Zunge $\Delta\Omega(1/2) \approx 5{,}6 \times
10^{-8}$, das Gesamtmaß aller Zungen $\approx 8 \times 10^{-4}$. Beides liegt
um fünf bis zwei Größenordnungen unter der fraktalen Korrektur \BM. Kein
Massenverhältnis wird durch Resonanz-Einrasten auf ein $p/q$ gezogen.

Daraus folgt eine Klarstellung, die im Korpus bisher nur implizit war: Die
exakten Rationalzahlen --- $43200$, $8^2 \cdot 5^2 \cdot 27$, die
Wicklungszahlen $(r_i, p_i)$ --- sind \emph{topologisch} (ganzzahlige
Wicklungen auf T$^4$), nicht \emph{dynamisch} (Einrasten). Diskretheit und
Kopplung sind in FFGFT zwei getrennte Mechanismen. Das ist konsistent mit
Dok.~328, wo $\varphi$ außerhalb aller Zungen liegt: Die Kopplung ist zu
schwach, um Verhältnisse festzuhalten; was festgehalten wird, hält die
Topologie.

% ============================================================
\section{Satz G$'$: $\xi$ als Auflösungsboden --- die Grenze ist stufenabhängig}
% ============================================================

\textbf{Frage.} Ist ausgeschlossen, dass $\xi$ bei der Resonanzfindung oder
Analyse eine Rolle spielt?

\textbf{Antwort.} Nicht absolut. Drei Rollen sind zu unterscheiden. Als
\emph{Auswahlkriterium} in einer Suchfunktion der Form
$\exp(-(\omega-\omega_0)^2/4\xi)$ ist $\xi$ ausgeschlossen: das Zentrum
$\omega_0$ tut die Arbeit, $\xi$ setzt nur die Breite \BM. Als
\emph{dynamische Kopplung} ist $\xi$ ausgeschlossen: $K_\mathrm{eff}=2\pi\xi$
ergibt Arnold-Zungen von $5{,}6\times10^{-8}$ (Satz~G$''$) \BM. Als
\emph{Auflösungsboden} hingegen ist $\xi$ die natürliche Skala: Dok.~328 setzt
$g/\omega = \xi$ als Linienbreite der Torusmoden, Dok.~162 die fraktale
Dämpfung pro Operation als $K_\mathrm{frak} = 1 - 100\xi$.

\subsection{Die 1,05\,\% sind eine fraktale Stufe}

Die Korrektur bare $\to$ gemessen aus Dok.~338 ($1{,}05\,\%$) und die
Dämpfungsstufe $100\xi = 1{,}33\,\%$ aus Dok.~162 stehen im Verhältnis $0{,}79$
\BM. Die Schwellen in Satz~F und~G wurden mit $\varepsilon = 1{,}05\,\%$
gerechnet --- sie sind damit die Schwellen \emph{einer} fraktalen Stufe, nicht
der Theorie. Wird diese Stufe analytisch abgezogen (sie ist nach Dok.~338/162
berechenbar), rückt die Trennschärfe auf $\xi$:

\begin{table}[H]
\centering
\caption{Schwellen aus Satz~F/G für die Auflösungsstufen des Korpus
  (\texttt{pruef\_343d\_xi\_aufloesung.py})}
\begin{tabular}{lrrrrl}
\toprule
Auflösung $\varepsilon$ & Wert & $Q_c$ & $p^*$ & Primen $\leq p^*$ & Galois trennscharf bis \\
\midrule
$1{,}05\,\%$ (Dok.~338) & $1{,}05\times10^{-2}$ & 12{,}5 & 9{,}8 & 4 & $k=3$ \\
$100\xi$ (Dok.~162) & $1{,}33\times10^{-2}$ & 11{,}1 & 8{,}7 & 4 & $k=3$ \\
$\xi$ (Dok.~328) & $1{,}33\times10^{-4}$ & 111 & 87 & 23 & $k\geq 8$ \\
$1/Q_\mathrm{TFLN}$ (Dok.~186) & $10^{-6}$ & 1283 & 1000 & 168 & $k\geq 8$ \\
$\xi^2$ & $1{,}78\times10^{-8}$ & 9619 & 7500 & 950 & $k\geq 8$ \\
\bottomrule
\end{tabular}
\end{table}

Der Galois-Raster liegt bei $k_\mathrm{max}=8$ noch bei $0{,}128\,\%$
mittlerem Abstand --- zehnmal gröber als $\xi = 0{,}0133\,\%$. Bei
$\xi$-Auflösung bleibt die Schicht mindestens bis $k=8$ trennscharf, und
Primzahlen bis $83$ hätten unterscheidbare Euler-Faktoren \BM.

\subsection{Photonik als Zugang zum $\xi$-Boden}

Dok.~186 gibt für TFLN-Resonatoren $Q > 10^6$. Zum Auflösen von $\xi$ genügt
$Q > 1/\xi = 7500$; zum Auflösen von $\xi^2$ wäre $Q > 5{,}6\times10^7$ nötig
\BM. Photonische Hardware reicht also bis auf den $\xi$-Boden, nicht darunter
--- während NISQ-Hardware nach Dok.~034 die $\xi$-Korrekturen im Rauschen
verliert. Das ist die einzige Stelle, an der $\xi$ ,,in der Analyse eine Rolle
spielt``: nicht als das Werkzeug, das die Resonanz findet, sondern als die
Skala, bis zu der überhaupt unterschieden werden kann. Dieselbe Rolle hat $\xi$
überall im Korpus: $43200 = (r_\mu/r_e)^2/\xi$ sagt, \emph{wo} die Resonanz
liegt, nicht wie man sie sucht.

\textbf{Vorbehalte \SM.} Ob die $100\xi$-Stufe tatsächlich abziehbar ist, ist
im Korpus behauptet, nicht durchgeführt. Dok.~186 gibt $Q>10^6$ für die optische
Phase, nicht für Massenverhältnisse; die Übertragung $1/Q \to \varepsilon$ ist
Analogie. Die Prüfung von Dok.~186 selbst (\texttt{pruef\_343e\_dok186\_check.py})
bestätigt die TFLN-Güte, widerlegt aber das dortige ,,Einrasten`` auf
$\varphi^{-k}$ (Widerspruch zu Dok.~328 [K]), die Fibonacci-Bedingung für
Wicklungszahlen ($r_\mu = 16/5$, $16 \notin F$) und die B-Meson-Deutung
($\alpha=\xi$ bräuchte $9\times10^8$ Ereignisse für $4\sigma$); Dok.~186 trägt
seit dem 3.~September 2026 einen entsprechenden Vorbemerkungskasten.

\subsection{Folge für die Faktorisierung: $N_\mathrm{max} \approx 1/\varepsilon$}

Ein resonanzbasierter Faktorisierer muss die Periode $r$ von $r{+}1$ trennen
(Auflösung $1/r$) bzw.\ nach der Präzisionsbarriere aus Dok.~075
($\Delta f \approx f_0/n^2$) den Faktor $n$ isolieren (Auflösung $1/n^2$).
Der Auflösungsboden $\varepsilon$ deckelt ihn daher bei $N_\mathrm{max}\approx 1/\varepsilon$ \BM:

\begin{table}[H]
\centering
\caption{Auflösungsboden und Deckel eines resonanzbasierten Faktorisierers}
\begin{tabular}{lrrrr}
\toprule
$\varepsilon$ & Faktoren bis $n_\mathrm{max}$ & Zahlen bis $N_\mathrm{max}$ & Bits \\
\midrule
$1{,}05\,\%$ (eine fraktale Stufe) & 10 & 95 & 7 \\
$\xi$ & 87 & $7{,}5\times10^{3}$ & 13 \\
$1/Q_\mathrm{TFLN}$ & $10^3$ & $10^{6}$ & 20 \\
$\xi^2$ & $7{,}5\times10^3$ & $5{,}6\times10^{7}$ & 26 \\
RSA-2048 & --- & bräuchte $\varepsilon\approx2^{-1024}$ & --- \\
\bottomrule
\end{tabular}
\end{table}

Das $p^*\approx 87$ aus der Schwellentabelle \emph{ist} die Faktorisierungsaussage:
Ein Resonanzverfahren auf dem $\xi$-Boden trennt Faktoren bis etwa $87$. Selbst
mit $\xi^2$ sind es 26~Bit --- Probeteilung erledigt das in Millisekunden.
Was sich mit dem $\xi$-Boden verschiebt, ist die Galois-Schicht der Physik
($k\approx3\to k\geq8$), nicht die Faktorisierung: Die Physik gewinnt
Trennschärfe bei Massenverhältnissen, die Zahlentheorie gewinnt nichts, weil
ihre Zahlen keine $\xi$-Skala tragen (Dok.~075).

\subsection{Warum klassische Verfahren weiter reichen: Zustand gegen Prozess}

Klassische Verfahren lösen nichts auf --- sie zählen. $a^r \bmod N = 1$ ist wahr
oder falsch; es gibt kein $\varepsilon$. Jedes Bit ist ein $\mathbb{Z}_2$-Zustand,
im Korpus eine Wicklungszahl auf T$^4$ (Dok.~247): ein \emph{Zustand}, der ohne
Energie existiert und nach jedem Rechenschritt neu auf $0$ oder $1$ geschwellt
wird. Das ist die topologische Diskretheit aus Satz~G$''$. Ein Resonanzverfahren
trägt seine Information dagegen als \emph{Prozess} --- Phase, Frequenz --- und
hat keinen Zustand, auf den es zurückschwellen könnte; der Boden $\varepsilon$
bleibt ihm erhalten. Der Preis der Exaktheit ist die Aufzählung:

\begin{table}[H]
\centering
\caption{Klassische Grenzen sind Zählgrenzen (Operationen), keine Auflösungsgrenzen}
\begin{tabular}{rrrrl}
\toprule
Bits & Probeteilung $\sqrt N$ & Pollard-$\rho$ $N^{1/4}$ & NFS & \\
\midrule
13 & 90 & 10 & $9\times10^{2}$ & Resonanzdeckel bei $\xi$ \\
26 & $8\times10^3$ & 90 & $3\times10^{4}$ & Resonanzdeckel bei $\xi^2$ \\
512 & $10^{77}$ & $3\times10^{38}$ & $2\times10^{19}$ & Stunden \\
2048 & --- & $3\times10^{154}$ & $1{,}5\times10^{35}$ & $5\times10^{12}$ Jahre bei $10^{15}$ op/s \\
\bottomrule
\end{tabular}
\end{table}

Zur Energie: Landauers Grenze $k_BT\ln 2$ gilt im Korpus unverändert und ohne
$\xi$-Modifikation (Dok.~184), aber \emph{nur für die Löschung} eines Bits
(Dok.~247); Dok.~255 leitet sie als eingefrorene Energie einer minimalen
Windungsdifferenz $\Delta w=1$ aus der T$^4$-Geometrie her. Sie deckelt die
Dissipation pro irreversiblem Schritt, nicht die Zahl der Rechenschritte;
reversible Schritte sind ausgenommen. Die klassische Wand ist deshalb die Zeit,
nicht die Energie \BM.

Die beiden Verfahrenstypen sind dual: \textbf{Exaktheit kostet Aufzählung,
Auflösung kostet Präzision.} Resonanz hat konstante Kosten pro Kandidat und
einen Boden bei $1/\varepsilon$; Zählung hat keinen Boden und Kosten, die
subexponentiell mit $N$ wachsen. Beides sind Grenzen; die Zählgrenze liegt
nur um Hunderte Größenordnungen weiter außen. Dok.~316 sagt es allgemein: Die
arithmetische Struktur lässt sich nicht in ein Kontinuum projizieren, ohne
Information zu verlieren.


\subsection{Quantenrechner: der empirische Maßstab}

Das Schaltkreismodell von Shors Algorithmus ist digital und von der
Weyl-Obstruktion (Satz~F) nicht betroffen: Die QFT läuft über $\mathbb{Z}_Q$,
$Q=2^{2n}$, mit gleichabständigen Eigenfrequenzen und linearer Zähldichte
$N(\lambda)\propto\lambda$; die Zahlentheorie steckt in der Periode $r$ von
$a^x \bmod N$, nicht im Spektrum \BM. Das gilt für das Modell.

Der physische Apparat ist ein Analogrechner mit quantenmechanischer Präzision.
Ein Qubit ist ein Zweiniveau-Resonator, jedes Gatter ein resonanter Antrieb,
ein Photon eine Mode mit kontinuierlicher Phase. Quantenfehlerkorrektur misst
nicht den Qubit-Zustand selbst, sondern Syndrome über viele Runden; der logische
Qubit-Zustand bleibt eine kontinuierliche Superposition, das Dekodieren ist ein
probabilistisches Schätzproblem auf dem Syndromgraphen. Das ist kein
Zurückschwellen auf Bit-Ebene, sondern Mittelung über viele Messungen ---
dasselbe Grundproblem wie beim klassischen Analogrechner: ob die Mittelung den
analogen Fehlerboden unter die Nutzgrenze drückt. Das Substrat ist gleichgültig
--- Photonik, Ionen, Supraleiter unterscheiden sich nur im Fehlerkanal. Die
photonischen Operationen in Dok.~186 (analoge Matrixfaltung, Resonator-Einrasten)
liegen mit $\varepsilon = 1/Q \approx 10^{-6}$ bei $N_\mathrm{max}\approx10^{6}$
(20~Bit, Satz~G$'$).

Die Frage ist dieselbe wie beim Analogrechner: ob die Mittelung skaliert. Der
Korpus hat keinen Mechanismus, der einen physischen Resonator von den
Resonanzgrenzen befreit (Dok.~034: $\xi$-Effekte unter dem Rauschen; Dok.~162:
$100\xi$-Dämpfung pro Gatter, Charakter offen). Die Behauptung des
Schwellentheorems ist \SM\ --- jede Bedingung einzeln demonstriert, das Produkt
bei Skalierung nicht. Entschieden wird es nicht in der Theorie, sondern auf der
Faktorisierungskurve.

\begin{table}[H]
\centering
\caption{Faktorisierung mit Shors Algorithmus gegen klassische Rekorde und Bausteine (Stand September 2026)}
\begin{tabular}{llll}
\toprule
Jahr & Shor auf Hardware & klassisch (NFS) & Bausteine \\
\midrule
2001 & $N=15$ & --- & --- \\
2012 & $N=21$ (mit Vorwissen) & --- & --- \\
2019 & Versuch $N=35$ gescheitert & --- & --- \\
2020 & --- & RSA-250 (829~Bit) & --- \\
2024 & --- & --- & Fehlerkorrektur unter Schwelle (Google Willow) \\
2025 & --- & --- & 48 logische Qubits (Quantinuum Helios) \\
2026 & --- & --- & 94 logische Qubits $<10^{-4}$; logische QFT (12 log.\ Qubits) \\
\bottomrule
\end{tabular}
\end{table}

Die größere ,,Quantenrekorde`` ($56\,153$, $4\,088\,459$, $2{,}6\times10^{14}$,
$1{,}1\times10^{12}$) betreffen Zahlen mit Spezialstruktur (Faktoren, die sich
in wenigen Bits unterscheiden), gelöst mit Annealing oder QAOA nach klassischer
Vorverarbeitung; Gutmann und Neuhaus (IACR ePrint 2025/1237) haben sie
sämtlich mit einem VIC-20 von 1981, einem Abakus und einem Hund repliziert.
Der größte \emph{simulierte} Shor-Lauf (39~Bit) lief klassisch auf 2048 GPUs.

Die Shor-Kurve steht seit 2012 bei $N=21$; die klassische Kurve bei 829~Bit.
Der Maßstab ist daher empirisch: \textbf{Ein Vorteil muss auf der
Shor-Faktorisierungskurve sichtbar werden, nicht auf Komponenten-Benchmarks.}
Das ist die Korpusposition aus Dok.~075 --- \textbf{[X]} kein allgemeiner
exponentieller Quantenvorteil nachgewiesen --- und sie ist im September 2026
unverändert gültig. Die theoretische Möglichkeit (Shor polynomial; Weyl greift nicht; ob eine
andere physikalische Obstruktion existiert, ist offen) bleibt \SM, keine [X];
ihre Beweislast liegt beim Experiment.
Auflösungsboden (G$'$), Zählgrenze (Zustand/Prozess) und Ressourcengrenze
(Shor) sind drei verschiedene Deckel; für die Faktorisierung wirkt heute
allein der dritte, und er ist bislang nicht durchbrochen.


% ============================================================
\section{Zusammenfassung}
% ============================================================

\begin{table}[H]
\centering
\caption{Ergebnisse und Status}
\begin{tabular}{lll}
\toprule
Satz & Inhalt & Status \\
\midrule
A & $\zeta_{T^4/\mathbb{Z}_3}(s) = (1-3^{-s})\,\zeta(s) = L(s,\chi_0)$ & \BM \\
A$'$ & $L(s,\chi_1)$ misst Twist-Sektor-Asymmetrie & \SM \\
B & Euler-Produkt hierarchisch nach $k=\mathrm{ord}_p(3)$ & \BM \\
B$'$ & Fraktale Dämpfung unterdrückt Euler-Faktoren $k\geq 7$ & \BM \\
C & Orbifold-Nullstellen bei $s = 2\pi ik/\ln 3$ auf $\mathrm{Re}(s)=0$ & \BM \\
C$'$ & Keine Resonanz Riemann-Nullstellen mit Z$_3$-Einheit & \BM \\
D & Galois-Charakter trennt nicht; Klassen $=(\mathbb{Z}/26)^*/\langle 3\rangle\cong\mathbb{Z}_4$ & \BM \\
D$'$ & Sektorpaarung $k\leftrightarrow -k$: $f_1\leftrightarrow f_2$, $f_3\leftrightarrow f_4$ --- 4 Klassen $\to$ 2 & \BM \\
D$''$ & Negation koppelt Ordnung 26 an 13; keine Sub-Leptonen-Schicht & \BM \\
D$'''$ & $\{f_1,f_2\}$ Majorana-Schicht, $\{f_3,f_4\}$ Dirac-Schicht (Sektorpaarung $=$ Majorana-Kriterium) & \BM \\
D$''''$ & $f_2=O_3$ aktive Neutrinos (Dok.~340); $f_1=O_1$ steril/rechtshändig, Seesaw-Partner & \SM \\
E & $p=13$ einzige Primzahl mit $\mathrm{ord}_p(3)=3$ weltweit; Primen durch $3^k-1$ erzwungen & \BM \\
E$'$ & $\mathrm{ord}_p(3)=k$, $3|k$ $\Rightarrow$ $p\equiv 1\pmod 3$ (Z$_3$-kompatibel) & \BM \\
F & Euler-Faktor-Schwelle $p^*\approx 9{,}76$ ($s=2$) = 7-limit-Grenze & \BM \\
F$'$ & Produktdichte: Zufallsniveau ab $p_\mathrm{max}\approx 30$ (±0{,}5\,\%) & \BM \\
F$''$ & Weyl-Obstruktion: Grenze liegt im Gegenstand, nicht im Werkzeug (Dok.~316) & \BM \\
G & Kopplung: $Q_c=\pi/\sqrt{6\varepsilon}\approx 12{,}5$ bei $1{,}05\,\%$; Grenze $k\approx3$--$4$ = kritischer Punkt $\kappa=1$ & \BM \\
G$'$ & Dreiteilung gleich, Gesetz verschieden: Arnold $K^q$ vs.\ Farey $1/q^2$ & \BM \\
G$''$ & FFGFT dynamisch tief unterkritisch ($K_\mathrm{eff}=2\pi\xi$); Diskretheit topologisch, nicht dynamisch & \BM \\
G$'''$ & Satz~F auf $\mathbb{Z}/N\mathbb{Z}$ (RSA) nur als Analogie & \SM \\
H & $1{,}05\,\%$ (Dok.~338) $\approx 100\xi$ (Dok.~162): Grenze $k\approx3$ ist die einer fraktalen Stufe & \BM \\
H$'$ & Bei $\varepsilon=\xi$: $Q_c\approx111$, $p^*\approx87$, Galois trennscharf bis $k\geq8$ & \BM \\
H$''$ & TFLN $Q>10^6$ löst $\xi$, nicht $\xi^2$; $\xi$ = Auflösungsboden, kein Auswahlkriterium & \BM \\
H$'''$ & Abziehbarkeit der $100\xi$-Stufe; Dok.~186: Einrasten/Fibonacci/B-Meson widerlegt & \SM \\
I & Resonanz-Faktorisierer gedeckelt bei $N_\mathrm{max}\approx1/\varepsilon$ (13--26 Bit); klassisch: Zählgrenze, Landauer nur Löschung & \BM \\
I$'$ & Schaltkreismodell digital (Weyl greift nicht); physischer Apparat = Resonator, unterliegt $\varepsilon$/Weyl bis Schwellentheorem demonstriert \SM; Kurve seit 2012 bei $N=21$ & [X] \\
\bottomrule
\end{tabular}
\end{table}

Die zentrale Aussage: Die volle Riemann-Zeta $\zeta(s)$ gehört zum flachen
$\mathbb{R}^4$-Spektrum. Die Zeta des T$^4$/Z$_3$-Orbifolds ist
$L(s,\chi_0) = (1-3^{-s})\,\zeta(s)$ — eine gefilterte Version, bei der
der 3er-Euler-Faktor herausgezogen und die Primzahlen nach der Galois-Hierarchie
$k = \mathrm{ord}_p(3)$ gewichtet sind. Das ist die Zeta-Funktions-Form
der Eintrittsregel aus Dok.~342.

% ============================================================
\section*{Prüfskripte}
% ============================================================

Verzeichnis \texttt{2/python/Dok343\_Skripte/}, Wrapper \texttt{run\_all\_343.sh}:
\begin{itemize}
  \item \texttt{pruef\_343\_zeta\_galois.py} ---
        Sätze~A--D; partielle Zeta-Zerlegung, L-Funktionen, Nullstellen,
        Resonanztest, gewichtetes Euler-Produkt.
  \item \texttt{pruef\_343b\_klassen\_symmetrie.py} ---
        Satz~D; Frobenius-Klassen in GF$(27)^*$ explizit, Wirkung von Inversion
        und Negation, Nebenklassen von $\langle 3\rangle$ in $(\mathbb{Z}/26)^*$,
        exakte Werte $L(s,\chi_1)$.
  \item \texttt{pruef\_344\_primzahlmuster.py} ---
        Sätze~E--F; Muster in $\mathrm{ord}_p(3)$, exakte Schwelle $p^*$,
        Trefferrate Galois vs.\ Zufall, Euler-Faktoren.
  \item \texttt{pruef\_343c\_kopplung\_zahlen.py} ---
        Abschnitt Klartext~II; Farey-Überdeckung und $Q_c(\varepsilon)$, Gesetzvergleich
        Arnold $K^q$ gegen Farey $1/q^2$, Regime von $K_\mathrm{eff}=2\pi\xi$,
        Galois-Überdeckung nach $p_\mathrm{max}$.
  \item \texttt{pruef\_343f\_zweite\_schicht.py} --- Satz~D$'''$; Klassen, Orbits,
        Involutionen $k\mapsto-k$ und $k\mapsto k^{-1}$, Dirac/Majorana-Zerlegung.
  \item \texttt{pruef\_343d\_xi\_aufloesung.py} --- Satz~G$'$; Schwellentabelle für
        $\varepsilon\in\{1{,}05\,\%,\,100\xi,\,\xi,\,1/Q_\mathrm{TFLN},\,\xi^2\}$.
  \item \texttt{pruef\_343e\_dok186\_check.py} --- Prüfung von Dok.~186: TFLN-Güte,
        Arnold-Zungen an Fibonacci-Näherungen, 50-GHz-Leiter, B-Meson-Statistik,
        Wicklungszahlen gegen Fibonacci.
\end{itemize}

Alle Assertions bestanden.

% ============================================================
\section*{Vermerk zur Verwendung von KI-Assistenz}
% ============================================================

Dieses Dokument wurde mit Unterstützung von Claude (Anthropic) erstellt.
Alle numerischen Berechnungen sind durch Python-Prüfskripte verifiziert.
Die physikalischen Interpretationen und Schlussfolgerungen stammen von
Johann Pascher (ORCID 0009-0000-6518-4064).

\begin{thebibliography}{9}
\bibitem{P057} J.~Pascher, \textit{Dok.~057: Relationales Zahlensystem}, FFGFT-Korpus.
\bibitem{P159} J.~Pascher, \textit{Dok.~159: Harmonische Struktur des Torus}, FFGFT-Korpus.
\bibitem{P034} J.~Pascher, \textit{Dok.~034: T0-Formalismus und QM-Optimierung}, FFGFT-Korpus.
\bibitem{P162} J.~Pascher, \textit{Dok.~162: T0-Optimierungsstrategien im Vergleich}, FFGFT-Korpus.
\bibitem{P186} J.~Pascher, \textit{Dok.~186: FFGFT-Photonik-Analyse}, FFGFT-Korpus (mit Vorbemerkung vom 3.~September 2026).
\bibitem{P184} J.~Pascher, \textit{Dok.~184: Das p-Bit in der FFGFT}, FFGFT-Korpus.
\bibitem{P247} J.~Pascher, \textit{Dok.~247: Information als Zustand und als Prozess}, FFGFT-Korpus.
\bibitem{P255} J.~Pascher, \textit{Dok.~255: Informationskinetik --- Landauer und Vopson aus der T$^4$-Geometrie}, FFGFT-Korpus.
\bibitem{P328} J.~Pascher, \textit{Dok.~328: Kopplungsregime, Resonanz und Synchronisation}, FFGFT-Korpus, August~2026.
\bibitem{P330} J.~Pascher, \textit{Dok.~330: Drei Operationen von T$^4$ zum Beobachtraum}, FFGFT-Korpus, August 2026.
\bibitem{P336} J.~Pascher, \textit{Dok.~336: Die algebraische Brücke: FFGFT und $\mathbb{Z}_3\mathbb{C}$}, FFGFT-Korpus, August 2026.
\bibitem{P075} J.~Pascher, \textit{Dok.~075: RSA und Periodenfindung}, FFGFT-Korpus, August~2026.
\bibitem{P147} J.~Pascher, \textit{Dok.~147: Quantencomputing und T0-Zeit-Masse-Dualität}, FFGFT-Korpus, August~2026.
\bibitem{P316} J.~Pascher, \textit{Dok.~316: FFGFT und die Riemannsche Zeta-Funktion}, FFGFT-Korpus, August~2026.
\bibitem{P342} J.~Pascher, \textit{Dok.~342: Faktorisierungsklassen in G$(6)$ und die Harmonik der Primzahlen}, FFGFT-Korpus, 2.~September 2026.
\bibitem{Gutmann2025} P.~Gutmann, S.~Neuhaus, \textit{Replication of Quantum Factorisation Records with an 8-bit Home Computer, an Abacus, and a Dog}, IACR ePrint 2025/1237.
\bibitem{Gidney2019} C.~Gidney, M.~Ekerå, \textit{How to factor 2048 bit RSA integers in 8 hours using 20 million noisy qubits}, arXiv:1905.09749; C.~Gidney (2025), Ressourcenschätzung $<10^6$ Qubits.
\bibitem{Dirichlet} P.\,G.\,L.~Dirichlet, \textit{Über die Bestimmung der Dichte einer unendlichen Anzahl von Punkten}, 1837.
\end{thebibliography}
